\documentclass[10pt,a4paper]{amsart}
\usepackage[margin=19mm]{geometry}
\usepackage{amsmath,amssymb,mathtools}
\usepackage[scr=rsfs,cal=euler]{mathalfa}
\usepackage{xfrac}
\usepackage[unicode,hidelinks,bookmarks=false]{hyperref}
\newcommand{\ie}{i.e.\ }
\newcommand{\cf}{cf.\ }
\newcommand{\resp}{respectively\ }
\newcommand{\Subsection}{\subsection}
\providecommand{\nobreakdash}{\nobreak\hbox{-}}
\setlength{\emergencystretch}{2em}
\multlinegap0pt
\usepackage{bm}
\usepackage{bbm}
\usepackage{dsfont}
\usepackage{xspace}
\usepackage{cancel}
\usepackage{mathtools}
\usepackage{amsthm}
\makeatletter
\@ifundefined{lemma}{\newtheorem{lemma}{Lemma}}{}
\@ifundefined{theorem}{\newtheorem{theorem}{Theorem}}{}
\makeatother
\title[Optimal Recovery for Solving Variational Problems]{Optimal Recovery for Solving Variational Problems}
\author{Ting Wang, Gideon Simpson, Jaroslaw Knap}
\date{Source: arXiv:2609.30049v1 (CC BY 4.0)}
\begin{document}
\maketitle

\noindent\emph{Source and licence.} Optimal Recovery for Solving Variational Problems. Version arXiv:2609.30049v1 (CC BY 4.0), distributed under the Creative Commons Attribution 4.0 International licence, as recorded in the arXiv licence metadata for this exact version and shown on the abstract page https://arxiv.org/abs/2609.30049v1. Full rights notice: licenses/arxiv-2609.30049-CC-BY-4.0.txt.\par\smallskip

\noindent\textit{Editorial scope.} Counted unit: the Theorem whose source label is \texttt{eqn:var\_JM} in the article (source0.tex lines 1704--1718, source label \texttt{eqn:var\_JM}), delivered together with the article's own complete original proof of it (source0.tex lines 1721--1805, one \texttt{proof} environment of 85 lines). No mathematics is written, repaired or re-derived: statement and proof are the article's own text line for line. Required context retained uncounted: eqn:var-problem-Sobolev (lines 703--706); eqn:uz (lines 743--747); eqn:u-star (lines 794--796); lem:quadratic-estimate (lines 1645--1653). Every definition, display and label of those lines is retained unchanged. Omitted from this package and not redistributed: 1--702, 707--742, 748--793, 797--1644, 1654--1703, 1719--1720, 1806--2992. Nothing inside a retained excerpt is omitted or altered: every retained body is delivered exactly as the article's lines are written, so no omission receipt of any kind accompanies this package. Citations: the retained excerpts carry no citation command at all, so no bibliography entry of the article is transcribed and no reference is redistributed. Reference adaptation: none was needed and none was made; every reference inside the delivered unit resolves inside it, so no unresolved reference or citation is produced. Printed slips kept literal: no printed word, symbol, hypothesis or number of the counted statement or of its proof was altered. Layout and class: the article's own class is \texttt{elsarticle} with packages including amsmath, amssymb, enumitem, mathrsfs, amsthm, graphicx, multicol, comment, xcolor, url, hyperref, dsfont, inputenc, siunitx; the wrapper is the lane's pinned \texttt{amsart} editorial wrapper at 10pt on A4, which restates the theorem-like environments the retained text needs and adds the article's own macro definitions that the retained text needs (none), with extra wrapper packages (bm, bbm, dsfont, xspace, cancel, mathtools). The delivered numbering is the unit's own. Assets: no retained span contains a figure environment, a graphics-inclusion command, a PDF-page inclusion, a table or a TikZ picture; no image file is copied into this package, the package declares no asset dependency and no image file is redistributed with it. Licence: version arXiv:2609.30049v1 of this article is distributed under the Creative Commons Attribution 4.0 International licence, as recorded in the arXiv licence metadata for this exact version and shown on the abstract page https://arxiv.org/abs/2609.30049v1. A full rights notice is delivered at licenses/arxiv-2609.30049-CC-BY-4.0.txt. Characters: every retained excerpt is pure ASCII, so the delivered engine, XeLaTeX with its default Latin Modern text font, reports no missing character.
\begin{equation}
  \label{eqn:var-problem-Sobolev}
  \inf_{u \in W^{1, p}(\Omega)} \left\{J(u) = \int_{\Omega} G(x, u, \nabla_x u)\, dx ~~\text{s.t.}~~ u = 0 ~\text{on}~\partial \Omega\right\}.
\end{equation}

\begin{equation}\label{eqn:uz}
  u_M^{\boldsymbol{z}}(x) = K(x, \boldsymbol{\phi})K(\boldsymbol{\phi}, \boldsymbol{\phi})^{-1}\begin{pmatrix}
    \boldsymbol{z}\\ 0
  \end{pmatrix},
\end{equation}

\begin{equation}\label{eqn:u-star}
u_M^* \triangleq u_M^{\boldsymbol{z}^*},
\end{equation}

\begin{lemma}\label{lem:quadratic-estimate}
Let 
$K: \bar{\Omega} \times \bar{\Omega} \to \mathbb{R}$ be
a positive definite kernel.
If $\{\boldsymbol{z}_n \}\subset \mathbb{R}^{M_{\Omega}}$ is a sequence with $\lim_{n \to \infty}|\boldsymbol{z}_n| = \infty$, then
$
\lim_{n \to \infty}\int_{\Omega} \left|\nabla u_M^{\boldsymbol{z}_n}(x)\right|^2 \, dx = \infty .
$
\end{lemma}

\begin{theorem}
Assume that
  $K: \bar{\Omega} \times \bar{\Omega} \to \mathbb{R}$ is a positive
  definite kernel whose RKHS $\mathcal{U}$ is continuously embedded in
  $W^{1, p}(\Omega)$. 
  Suppose that  $\mathcal{U}$ contains the unique minimizer $u^*$ to the
    variational problem~\eqref{eqn:var-problem-Sobolev}.
  Then, for each fixed $M \in \mathbb{N}_+$, the
  minimization problem
  \begin{equation}
    \label{eqn:var_JM}
    \inf_{u \in \mathcal{U}_M} F_M(u)
  \end{equation}
  admits a minimizer $u_M^*$ defined by~\eqref{eqn:u-star}.
\end{theorem}

\begin{proof}
Recall $u_M^{\boldsymbol{z}}$ as defined in~\eqref{eqn:uz}
and the function $I_M: \mathbb{R}^{M_{\Omega}} \to \mathbb{R}$ defined by
  \[
    I_M(\boldsymbol{z}) \triangleq
    J\left(u_M^{\boldsymbol{z}}\right) = F_M\left(u_M^{\boldsymbol{z}}\right),
  \]
  where the last equality holds since
  $u_M^{\boldsymbol{z}} \in \mathcal{U}_M$ and
  $u_M^{\boldsymbol{z}}(x_i) = 0 ~\text{for}~ i = M_{\Omega}+1,\ldots,
  M$.  
We shall show that $I_M$ is lower semicontinuous and there exists a non-empty compact sublevel set $\mathcal{S}$ so that $I_M$ attains its infimum at some $u_M^* \in \mathcal{S}$.  
  
  
{\bf Lower semicontinuity.}  Let $\{\boldsymbol{z}_n\}$ be a
  sequence converging to $z$ in $\mathbb{R}^{M_{\Omega}}$. To
  ascertain lower semicontinuity we shall show that
  $\liminf_n I_M(\boldsymbol{z}_n) \geq I_M(\boldsymbol{z})$. 
  
  
  
  Assume that
  $\lim_n |\boldsymbol{z}_n| = |\boldsymbol{z}| < \infty$, then the sequence
  $\{\boldsymbol{z}_n\}$ is bounded (in $n$). Thus, the sequence of
  RKHS norms
   \[\left\{\left\|u_M^{\boldsymbol{z}_n}\right\|_{\mathcal{U}}^2\right\} = \left\{\boldsymbol{z}_n^{\text{T}} K^{-1}(\boldsymbol{x}, \boldsymbol{x})  \boldsymbol{z}_n\right\}\] 
   is a convergent sequence with a finite limit and hence it is bounded in $n$.  Since $u_M^{\boldsymbol{z}}$ is linear in $\boldsymbol{z}$,
  it is immediate that
  $u_M^{\boldsymbol{z}_n}(x) \to u_M^{\boldsymbol{z}}(x)$ for all
  $x \in \Omega$. The boundedness of
  $\left\{\|u_M^{\boldsymbol{z}_n}\|_{\mathcal{U}}\right\}$ together with
  pointwise convergence leads to the weak convergence
  $u_M^{\boldsymbol{z}_n} \rightharpoonup u_M^{\boldsymbol{z}}$ in
  $\mathcal{U}$ as $n \to \infty$.  Recall that we assume that  $\mathcal{U}$ is
  continuously embedded in $W^{1,p}(\Omega)$, we further have
  $u_M^{\boldsymbol{z}_n} \rightharpoonup u_M^{\boldsymbol{z}}$ in
  $W^{1,p}(\Omega)$ as $n \to \infty$.  Applying the
  weakly lower semi-continuity of $J$ leads to the liminf property
\[
  \liminf_n I_M({\boldsymbol{z}_n}) = \liminf_n
  J\left(u_M^{\boldsymbol{z}_n}\right) \geq J\left(u_M^{\boldsymbol{z}}\right) =
  I_M({\boldsymbol{z}}).
\]
  
  
  
%%  
%  It remains to consider the case when $\lim_n |\boldsymbol{z}_n| = \infty$, then
%  by the coercivity of $I_M$ we immediately have
%  $\lim_n I_M(\boldsymbol{z}_n) = \infty$ and hence the liminf
%  property holds trivially.  
  
  
  

 


{\bf Coercivity.}
We first show that $\inf_{\boldsymbol{z}} I_M(\boldsymbol{z})$ has an upper bound.
Let $u^* \in \mathcal{U}$ be the minimizer of $J$.
Denote $\boldsymbol{z}^{\dagger} = u^*(\boldsymbol{x}_{\Omega}) \in \mathbb{R}^{M_{\Omega}}$, that is, the vector obtained by evaluating $u^*$ at the interior collocation points. 
Note that, by definition, $u_M^{\boldsymbol{z}^{\dagger}} \in \mathcal{U}_M$ and $u_M^{\boldsymbol{z}^{\dagger}}(x_i)=0$ for $i = M_{\Omega} + 1, \ldots, M$.
Hence, we obtain 
\[
\inf_{\boldsymbol{z} \in \mathbb{R}^{M_{\Omega}}}I_M(\boldsymbol{z}) = \inf_{u_M^{\boldsymbol{z}} \in \mathcal{U}_M} F_M(u_M^{\boldsymbol{z}}) \leq J(u^{\boldsymbol{z}^{\dagger}}) < \infty.
\]
Therefore, the minimization of $I_M(\boldsymbol{z})$ can be restricted to the sublevel set of $I_M$, defined by 
\[
\mathcal{S} = \left\{\boldsymbol{z} \in \mathbb{R}^{M_{\Omega}}:  I_M(\boldsymbol{z}) \leq J\left(u^{\boldsymbol{z}^{\dagger}}\right)\right\}.
\]
Clearly, $\mathcal{S}$ is non-empty since $\boldsymbol{z} ^{\dagger}\in \mathcal{S}$.
We further claim that $\mathcal{S}$ is compact. 
Since $I_M$ is lower semicontinuous with respect to $\boldsymbol{z}$, the set $\mathcal{S}$ is automatically closed. 
To see $\mathcal{S}$ is bounded, we prove by contradiction. Suppose $\mathcal{S}$ is unbounded, then there exists a sequence $\boldsymbol{z}_n \in \mathcal{S}$ with
$|\boldsymbol{z}_n| \to \infty$. By Jensen's inequality, this implies that
\[
I_M(\boldsymbol{z}_n) \geq \int \alpha_1 \left|\nabla u_M^{\boldsymbol{z}_n}(x)\right|^p \, dx - |\alpha_2| |\Omega| 
\geq \alpha_3 \left(\int \left|\nabla u_M^{\boldsymbol{z}_n}(x)\right|^2 \, dx\right)^{\frac{p}{2}} - |\alpha_2| |\Omega| 
\]
for some $\alpha_3>0$. 
Taking $n \to \infty$ leads to $I_M(\boldsymbol{z}_n) \to \infty$ by Lemma~\ref{lem:quadratic-estimate}, contradicting to the definition of $\mathcal{S}$!
This proves that the set  $\mathcal{S}$ is compact.
Therefore, $I_M$ attains a minimum $u_M^* \in \mathcal{S}$ since a lower semi-continuous function attains a minimum over a compact set.  
\end{proof}

\end{document}
